% Effective parameters, R0, identifiability --- plan \S4.11 (Phase 3).
% Sources: new_notes3 Part D \S20; new_notes Part II \S10--\S12;
% new_notes2 \S8. The p_eff example values and the trade-off inputs were
% re-verified numerically on 2026-08-07.

\section{Effective parameters, $R_0$, and identifiability}
\label{sec:peff-section}

The renewal model \eqref{eq:Icell}--\eqref{eq:Vfree} does not merely
disagree with classical BMVR in transients; it contains classical BMVR as an
exact special case, and the relationship between the two is precise enough
to be turned into predictions.

\subsection{Exact reduction in the exponential phase}

Suppose the incidence grows exponentially, $i(t)=i_0\ee^{rt}$, and look at
times late enough that the lower limits of the convolutions are irrelevant.
Then
\begin{equation}
\Icell(t)=i_0\ee^{rt}\Lap{\Ihat}(r),
\qquad
\Vfree'(t)=i_0\ee^{rt}\,\delta\Lap{K}(r)-c\,\Vfree(t),
\end{equation}
where $\Lap{f}(r)=\int_0^\infty\ee^{-r\alpha}f(\alpha)\,\dd\alpha$. A
classical model $\Icell'=i-d_{\Icell}\Icell$, $\Vfree'=p\Icell-c\Vfree$
driven by the same incidence has $\Icell(t)=i_0\ee^{rt}/(r+d_{\Icell})$ and
the eigenvalue relation $(r+d_{\Icell})(r+c)=\gamma Tp$. Matching the two
exponential solutions fixes
\begin{equation}
p_{\rm eff}(r)=\frac{\delta\,\Lap{K}(r)}{\Lap{\Ihat}(r)},
\qquad
d_{\Icell,{\rm eff}}(r)=\frac{1}{\Lap{\Ihat}(r)}-r.
\label{eq:peff}
\end{equation}

This reduction is \emph{exact}, not approximate: substituting
\eqref{eq:peff} into $(r+d_{\Icell})(r+c)=\gamma Tp$ returns
\begin{equation}
r+c=\gamma T\,\delta\Lap{K}(r),
\label{eq:char}
\end{equation}
which is precisely the characteristic equation of the renewal system. Thus
for every growth rate $r$ there is a unique constant-parameter BMVR that
reproduces the renewal model's exponential solution exactly. What
\eqref{eq:peff} cannot do is reproduce transients: a single pair
$(p,d_{\Icell})$ matches one growth rate, and a transient samples a
continuum of them (Figure~\ref{fig:D-exp} shows the exponential case; the
overlays of \S\ref{sec:overlays} show everything else).

The two ends of the map $r\mapsto(p_{\rm eff},d_{\Icell,{\rm eff}})$ are
explicit. Using $\int_0^\infty\delta K=V_\infty$ (from $V'=\delta K$) and
$\int_0^\infty\Ihat=\mean{T_{\rm prod}}$ (Proposition~\ref{prop:lifetime}),
\begin{equation}
r=0:\qquad
p_{\rm eff}(0)=\frac{V_\infty}{\mean{T_{\rm prod}}}
=\frac{\beta\,V_\infty}{\log\bigl(a/(a-1)\bigr)},
\qquad
d_{\Icell,{\rm eff}}(0)=\frac{\beta}{\log\bigl(a/(a-1)\bigr)};
\label{eq:peff0}
\end{equation}
the $r=0$ value of $p_{\rm eff}$ coincides with naive moment-matching ---
lifetime output divided by lifetime --- a pleasant consistency check, and
the matching used in the overlays of \S\ref{sec:overlays}. At the other
end,
\begin{equation}
r\to\infty:\qquad
p_{\rm eff}(r)\to\delta,
\qquad
d_{\Icell,{\rm eff}}(r)\to\mu+\delta,
\label{eq:peffinf}
\end{equation}
the \emph{young-cell} limit: in a fast-growing epidemic almost every
infected cell was infected moments ago, holds one particle, bursts at rate
$\delta$, and releases one particle when it does. Between the two ends,
$p_{\rm eff}$ decreases monotonically in $r$ across every parameter set we
examined (Figure~\ref{fig:peff-curves}): fast epidemics are dominated by
young cells and release slowly per cell; slow or stationary ones are
dominated by old cells and release quickly. At $(\beta,\mu,\delta)=(1,0,0.1)$:
$p_{\rm eff}(0)\approx4.59$ against $p_{\rm eff}(\infty)=0.1$ --- the same
pathogen, two orders of magnitude apart, depending only on how fast the
epidemic around it is growing.

\begin{figure}[H]
\centering
\includegraphics[width=0.92\textwidth]{figures/peff_dr_curves.pdf}
\caption{The effective release and removal rates \eqref{eq:peff} as
functions of the epidemic growth rate $r$, for three parameter sets.
Dotted and dashed lines mark the limits \eqref{eq:peff0} and
\eqref{eq:peffinf}. $p$ is not a constant of the pathogen; it is a
function of the epidemic's phase.}
\label{fig:peff-curves}
\end{figure}

The closed hypergeometric forms of $\Lap{\Ihat}$ and $\delta\Lap{K}$ are
given in Appendix~\ref{app:transforms}; the curves above are computed by
direct quadrature.

\subsection{The characteristic equation and the invariance of $R_0$}

Setting $r=0$ in \eqref{eq:char} gives the threshold condition
$\gamma T\,\delta\Lap{K}(0)=c$, i.e.\ $\gamma T V_\infty=c$, so the
basic reproduction number of the renewal model is
\begin{equation}
R_0^{\mathrm{ODE}}=\frac{\gamma T\,V_\infty}{c}.
\label{eq:R0ODE}
\end{equation}

\begin{theorem}[$R_0$ is unchanged by bursting]
\label{thm:R0inv}
At fixed total per-cell output $V_\infty$, replacing constant budding by
bursting leaves $R_0$ exactly unchanged. Bursting alters the
generation-time distribution --- hence the growth rate $r$, hence the
extinction probability, hence the stochastic early phase --- but not the
threshold itself.
\end{theorem}

This is the classical $R_0$-versus-$r$ distinction in a setting where both
sides are computable in closed form, and it connects the renewal model
directly to the ``perilous beginning'' theme of Chapter~3 and to the
flooding question of \S\ref{sec:flooding}. Two cautions accompany it.
First, the numerical checks of Figures~\ref{fig:E-growth} and
\ref{fig:F-R0}: the late growth rate of \eqref{eq:Icell}--\eqref{eq:Vfree}
agrees with the characteristic root $r_*$ of \eqref{eq:char} to
$\sim10^{-7}$ across five scenarios, and trajectories decay for
$R_0<1$ and grow for $R_0>1$, the transition resolved already at
$R_0=1.05$. Second, a notational care: \eqref{eq:R0ODE} is the
\emph{continuum} reproduction number, in which one virion contributes mean
$\gamma T/c$ new infections. A discrete branching process in which virions
are absorbed on infection carries instead the mean
$m=qV_\infty$, $q=\gamma T/(\gamma T+c)$; the two agree when
$\gamma T\ll c$, and it is $m$ that governs the branching calculations of
\S\ref{sec:flooding}. We write $R_0$ carefully in each context.

\begin{figure}[H]
\centering
\begin{subfigure}{0.47\textwidth}
\centering
\includegraphics[width=\textwidth]{figures/E_growth_rate_match.pdf}
\caption{Late growth = characteristic root.}
\label{fig:E-growth}
\end{subfigure}
\hspace{\fill}
\begin{subfigure}{0.47\textwidth}
\centering
\includegraphics[width=\textwidth]{figures/F_R0_threshold.pdf}
\caption{Sharp threshold at $R_0=1$.}
\label{fig:F-R0}
\end{subfigure}
\caption{Two structural checks of the renewal system: the fitted late growth
rate of $\Vfree$ against the root of \eqref{eq:char} (five scenarios,
relative error $\sim10^{-6}$--$10^{-7}$), and the $R_0$ threshold scan
(eight values through $R_0=1$).}
\label{fig:EF}
\end{figure}

\subsection{Identifiability: a negative result}
\label{sec:identifiability}

How much of the intracellular process can a population-level fit recover?
At any fixed growth rate $r$, the map
\begin{equation}
(\beta,\mu,\delta)\ \longmapsto\ \lrb{p_{\rm eff}(r),\,d_{\Icell,{\rm eff}}(r)}
\end{equation}
sends three rates to two numbers, and is therefore not invertible. This
closes, negatively, the ambition of the original Chapter~4 attempt, which
hoped that a good estimate of $p$ would grant access to the rates
constituting the quasi-stationary mean: it would not, and neither would
$(p,d_{\Icell})$ together. A third observable is required. In order of
experimental accessibility, the natural candidates are:
\begin{enumerate}
\item the burst-size dispersion --- the geometric ratio $1/a$ fixes one
combination of the rates directly (\S\ref{sec:burst-size});
\item the fraction of cells that clear without bursting, which is $b$;
\item the shape --- not merely the mean --- of the release-time density
$\varphi=\delta J$ (\S\ref{sec:burst-time}).
\end{enumerate}
Any one of these, together with $(p,d_{\Icell})$, over-determines
$(\beta,\mu,\delta)$ and converts the fit into a consistency test of the
whole intracellular model. The negative result is worth stating plainly: it
is not a limitation of a fitting procedure but of the information content
of constant-parameter BMVR itself.

\subsection{Prediction for fitting practice}

Since $p_{\rm eff}$ depends on $r$ and decreases with it
(Figure~\ref{fig:peff-curves}), BMVR fits to acute-phase data --- fast
growth, young cells --- and to set-point data --- slow or stationary ---
should yield systematically different estimates of $p$, in a direction and
by an amount the theory specifies. Equivalently: the eclipse-like delay
that within-host models usually insert by hand emerges here from the
intracellular dynamics with no extra parameter at all --- the renewal
system's early lag, visible in Figures~\ref{fig:overlay-V} and
\ref{fig:growth-phase}, \emph{is} the eclipse.
